\paragraph*{The published radiomics pipeline.} We ran the published MU radiomics pipeline \citep{christodoulou2026}
with the authors' public code and settings, on the same patients, masks, label and folds as PRISM. It extracts 4,892
radiomic features from the original, Laplacian of Gaussian and wavelet images of the four sequences of the current
scan, keeps 256 of them within each training set by the authors' ranking, and fits a class weighted gradient boosted
model tuned over 400 Bayesian search trials with patient grouped inner folds, followed by Platt calibration. The
authors' ranking adds permutation importance only when an inner fold holds at least 30 scans, so with our 13 to 16
scans it ranked features by coefficient size alone. We left the code unchanged. The published study gave every scan
of a patient who ever progressed the label of progression. We kept PRISM's label, one scan per documented
progression, so the comparison is of methods on one task, not a replication of the published estimate. The plan,
committed before the run, stated that the comparison with PRISM's imaging model would be withdrawn if the pipeline
did not rank below it.

% Sources: MASK_SENSITIVITY_PLAN.md; METHODS_RESULTS.md E5; SEGMENTER_STAGE_PLAN.md; SEGMENTER_STAGE.md;
% SEGMENTER_BRATS24_PLAN.md; CHRISTODOULOU_BRATS_MASKS_PLAN.md.
\paragraph*{Segmentation comparison.} We measured how each segmenter changes the results, in three planned analyses.
First, on MU we replaced every supplied mask with the Swin UNETR label map, kept the resection cavity of the supplied
mask, and refitted the imaging and combined models, the time-and-enhancement model and the published radiomics
pipeline in the same folds. Second, we ran PRISM from Swin UNETR alone at all three centres, with its 22 measurements
and nothing taken from the supplied masks, refitting on MU and applying the models unchanged externally. This plan was
written before any segmenter-derived number was linked to a label at LUMIERE or UCSF-ALPTDG. At UCSF-ALPTDG the
time-and-enhancement model was applied here with the approximated window. Third, the BraTS 2024 model replaces Swin
UNETR in the second analysis, with all 27 measurements. Each analysis compares each score with its value on the
supplied masks by paired AUC differences, and the rule asked that every paired CI include zero or lie above it, and
that the reading of the combined model's gain over the time-and-enhancement model be the same under both mask sets.
We also report agreement with the supplied masks (Dice) and the repeatability of the enhancing volume, as the median
fold change between consecutive scans at most 30 days apart, for which a single segmenter applied to every scan was
expected to be more stable.

% Sources: NUMBERS.md sections 6, 9, 13, 14; METHODS_PLAN.md; CLINICIAN_BASELINE_PLAN.md.
\paragraph*{Tests fixed in advance.} Table~\ref{tab:2} lists each planned test, the rule fixed in advance and the
outcome, including the two tests that passed only after a revision. The primary results are the external AUCs, the
gain over the time-and-enhancement model and the comparison with the published radiomics pipeline. The other tests
answered separate questions, and we made no correction for multiple tests. Results of analyses that were not fixed in
advance are called exploratory where they appear.

% Sources: NUMBERS.md section 9; METHODS_RESULTS.md; CHRISTODOULOU_REPLICATION.md; SEGMENTER_STAGE.md.
\begin{table}[!htbp]
\centering
\caption[Planned tests]{\textbf{Planned tests and their outcome.} Each test, the rule committed to the repository
before it, and the outcome. AUC differences with 95\% CIs.}
\label{tab:2}
\begin{threeparttable}
\footnotesize
\zebrarows
\begin{tabularx}{\linewidth}{@{}R{4.0cm}R{3.9cm}LR{0.9cm}@{}}
\toprule
Test & Rule fixed in advance & Outcome & Section \\
\midrule
Published radiomics pipeline against the clinical-history model on MU folds & Estimate only & $-$0.24 ($-$0.34 to $-$0.15) & \ref{sec:3.1} \\
Published radiomics pipeline below the imaging model on MU folds & Otherwise the comparison is withdrawn & \textbf{Yes}. $-$0.13 ($-$0.23 to $-$0.02) & \ref{sec:3.1} \\
MU results unchanged with Swin UNETR masks (supplied cavity kept) & Every paired CI includes 0, time-and-enhancement reading unchanged & \textbf{No} for the imaging model's AUC. \textbf{Yes} for the combined model, the time-and-enhancement model and the reading & \ref{sec:3.3} \\
Published radiomics pipeline below the imaging model on the Swin UNETR masks & Otherwise the claim is limited to the supplied masks & \textbf{Yes} & \ref{sec:3.1} \\
Results at three centres unchanged with PRISM's Swin UNETR stage & Every paired CI includes 0 or lies above it, gain reading unchanged at MU and LUMIERE & \textbf{No} for the MU imaging model. \textbf{Yes} for the other eight scores and the reading & \ref{sec:3.3} \\
Results at three centres unchanged with the BraTS 2024 stage & As for Swin UNETR & \todo{S4 PENDING} & \ref{sec:3.3} \\
Image measurements and schedule inputs add to the clinical-history model on MU & Gain $\geq$ 0.03, CI above 0 & \textbf{No}. $+$0.02 ($-$0.03 to $+$0.06) & \ref{sec:3.2} \\
MU AUCs unchanged with all scans before progression as negatives & Each changes by $<$ 0.05, time-and-enhancement reading unchanged & \textbf{No} for the imaging model, $-$0.054. \textbf{Yes} for the others and the reading & \ref{sec:3.1} \\
Unknown IDH and MGMT status as its own category & Every paired CI includes 0 & \textbf{Yes} & \ref{sec:2.2} \\
Deep networks tuned on MU add to the clinical-history model with image measurements & Gain $\geq$ 0.03, CI above 0 & \textbf{No}. No network gained & \ref{sec:3.2} \\
Combined model at LUMIERE & AUC $\geq$ 0.65 & \textbf{Yes} & \ref{sec:3.1} \\
Imaging model at UCSF & AUC $\geq$ 0.65, lower bound $>$ 0.5 & \textbf{Yes} & \ref{sec:3.1} \\
Disagreement of the two models marks scans the combined score ranks less well & AUC difference, CI above 0 & \textbf{Yes} at MU; not at LUMIERE or UCSF & \ref{sec:3.4} \\
Combined model over the time-and-enhancement model & CI above 0 & \textbf{No clear gain} at MU or LUMIERE & \ref{sec:3.2} \\
Marginal conformal sets, combined model & Coverage $\geq$ 0.78 in each cohort & \textbf{Yes} on point estimates & \ref{sec:3.4} \\
Reports pass the checker & $\geq$ 95\% in each cohort & \textbf{No} as planned (every LUMIERE report failed). \textbf{Yes} after one revision of the record fields and the prompt & \ref{sec:3.4} \\
Verifier closes a blind spot & $\geq$ 95\% of synthetic errors of that kind rejected & After one revision of the rules, \textbf{Yes} for direction and swapped volumes, \textbf{No} for a swapped tier word & \ref{sec:3.4} \\
Usefulness to clinicians & Not tested & Future work & \ref{sec:4} \\
\bottomrule
\end{tabularx}
\begin{tablenotes}\footnotesize
\item Each plan and the date it was committed are listed in Supplementary Section~S5. The conformal sets of the imaging model at LUMIERE (0.71) and the clinical-history
model at UCSF (0.70) did not meet the coverage rule.
\end{tablenotes}
\end{threeparttable}
\end{table}

% Sources: NUMBERS.md 2.3; SYSTEM_RESULTS.md; METHODS_RESULTS.md (bootstrap seeds).
\paragraph*{Statistics.} AUCs carry 95\% confidence intervals (CIs) from 1,000 bootstrap samples of patients (2,000
for the time-and-enhancement model, the AUC split, the group-conditional sets, the segmentation comparisons, the
published pipeline on the Swin UNETR masks, and the calibration, spectrum and marker-coding analyses). The bootstrap
resamples the held-out predictions without refitting the models, and paired differences use the same samples for
both scores. Calibration is summarised by the calibration slope and intercept.

% =====================================================================================================
